\chapter{Discussion}

\section{Summary of Findings}

Of the three retrieval-design factors, indexing representation moved retrieval
quality substantially, while chunk size produced no detectable difference and
retrieval strategy none without enrichment. The enrichment effect was not
uniform: it held for dense and hybrid retrieval at both chunk sizes on both
metrics but for BM25 only at the smaller chunk size and only on MRR@5, and
across question types it was present among factual questions, each naming one
company. Downstream, the better retrieval produced no detectable change in the
reported arm's answer behavior.

\section{Theoretical Contributions}

\subsection{Retrieval Design Is Dominated by What Enters the Index}

Chunk size produced no detectable difference, and neither did retrieval
strategy without enrichment. These are the two choices a retrieval design
usually turns on, and both nulls trace to one cause: what entered the indexed
string. Each chunk carried its transcript's company and period as metadata, but
by construction neither entered the text-only indexed string (Methods). Across
\thesisNumCompanies{} banks' near-identical calls over \thesisNumQuarters{}
quarters, identity reached the retrievers only where a transcript's own words
named it. It was missing from the search space rather than under-weighted by
the ranker, so neither chunk size nor retrieval strategy, as varied here, could
supply it. Where documents
differ in subject matter, content terms already discriminate and the effect
should be expected to shrink.

Results places the effect where this mechanism predicts: in factual questions,
each naming one company, and not in thematic questions, which name none. The
thematic null is a floor effect, consistent with the mechanism rather than
confirmation of it, and no evidence that identity metadata is irrelevant to
thematic retrieval.

Enrichment also displaced anchors: some questions retrieved one under text-only and
none under any enriched condition. They were marginal in
the baseline, most scoring in only one of the six text-only conditions, and their
anchors' chunks were still ranked in the enriched index, out of the top five;
for the anchors that moved, truncation does not explain the loss.
Identity in the prefix competes with content terms for the same top five, so
enrichment reallocates ranking capacity toward identity and away from content.
It is no free improvement: net positive here, uneven, with BM25 gaining least
(Figure~\ref{fig:coverage}), and with losers.

H3 corroborates the improvement \textcite[9]{yousuf2026} report on a corpus
outside the structured collections to which they expect it to carry
\parencite[13]{yousuf2026}: spoken transcripts that follow no template. Here
enrichment was crossed with chunk size and retrieval strategy, and followed
through to the answers. The split BM25 verdict and the displaced questions
qualify that corroboration without contradicting it.

\subsection{Retrieval Gains Did Not Reach the Answers}

Retrieval improved substantially with enrichment, and the reported arm's answer
outcomes did not follow: no difference in outcomes between conditions was
detectable. As with the retrieval nulls above, that failure to detect a
difference, at this sample size, is not evidence that the conditions answered
alike.

Two readings fit the result, and the design does not choose between them. The
generator may already have been handling what it was given, the additional
retrieved anchors falling in questions it would have abstained on either way.
Or the outcome classification, a four-way mechanical label, may be too coarse
to register the change. The second is the stronger reading, because the
instrument's limits are visible in the comparison itself. Over the six text-only conditions
only 65 of 120 pairwise comparisons were testable, the rest having identical
outcomes on every matched question, and within Case C each of those pairs
matched very few questions. The instrument, not the phenomenon, may be the
binding constraint.

Part of what is counted as a failure to convert may not be one. Case C
membership turns on any single anchor, so some items counted as unconverted were
ones where abstention was the correct response, which the next subsection takes
up.

\subsection{Partial Context Must Be Read Against the Question's Evidence
Requirement}

This is a contribution to benchmark design and evaluation: nothing in the
argument depends on which generator ran. Question
types differ in how much of their evidence they require. A thematic question is
disjunctive, any subset of its anchors supporting an answer, whereas a
comparative question is conjunctive, needing the anchors of both its companies.
Case C membership, however, turns on any single anchor, so a comparative item
holding one of its two anchors is in Case C although what was retrieved cannot
support the comparison. Abstention on such an item is the correct response, and
the scheme records it as a failure to convert. The error lies in the membership
rule, not in the generator and not in the answer.

The coverage split shows the difference is observed: at
partial coverage, disjunctive thematic observations mostly converted and
conjunctive comparative ones mostly abstained. The breakdown is descriptive,
computed, and unchecked. It does not make the downstream null an artifact; it
shows that some of what was counted as non-conversion was correct behavior,
without saying how much.

This is where the information processing account set out in Theory applies:
effectiveness depends on a match between what a task requires and what is
supplied \parencite[619, 622]{tushman1978}, and partial coverage met the
disjunctive requirement and fell short of the conjunctive one. The reading concerns the questions' requirements alone; it does
not test the theory, and no step is taken from a retrieval-design change to a
change in information processing capacity.

A partial-context threshold therefore has to depend on the question's evidence
requirement. Most benchmarks leave that requirement unrecorded, so even an
evaluator who wants to cannot apply the distinction. This is the gap: a
benchmark should carry the evidence requirement as a field, the way this one
carries question category.

\subsection{Automatic Grounding Checks Mislabel Faithful Text}

Every record the reported arm labeled ungrounded was read, because that stratum
was taken whole, so the finding is a census: the one claim of the manual review
resting on a full population. Of those
\thesisQeightAnsweredUngrounded{} records, four, nearly half, read as faithful.
Their labels trace to instrument findings I1 to I3. A year read correctly from
the context is verified as though it were a quantity; a figure the answer
computes and marks as computed has nothing verbatim to match; and a faithful
paraphrase shares no verbatim run with its source long enough to match.

These misreadings run in one direction: a verbatim-substring grounding proxy
penalizes answers that paraphrase, normalize, or derive, which are the behaviors
a good answer exhibits. An evaluation using
such a proxy measures copying as much as faithfulness, and should report the
direction of that bias rather than presenting the proxy's output as a
faithfulness rate.

A narrower point follows from I4: a thematic question about banks generally can
be answered accurately, well attributed, and fully responsive, while carrying the
off-target label. Those records came from a stratified draw, so they are no
estimate for the arm.

\section{Practical Contributions}

\subsection{For Systems Retrieving Over Near-Identical Documents}

The audience is teams retrieving over document sets whose members are
structurally alike and differ mainly in entity and period: earnings calls,
filings, contracts, incident reports, clinical notes. They should put the
identity metadata they already hold into the indexed string. The cost is a
short fixed prefix on every chunk, small beside re-chunking a corpus or
standing up a second retriever, where tuning effort usually goes.

The advice holds under three conditions: the metadata must discriminate, since
where documents differ in subject matter it can be dead weight; every prefixed
token consumes the encoder's content budget; and identity competes with content
for ranking capacity, so expect named losers as well as gains.

Expect the gain where a question names one company. Where questions name none,
both arms sat near floor here, so no difference was detectable either way, and
where a question needs both sides of a comparison, retrieving one side leaves it
unanswerable, which is a property of the question rather than of the index.

The token budget is an engineering trap: a prefix pushes the larger chunks past
the encoder's limit, and without a cut the encoder would drop the tail while the
lexical retriever indexed the whole string, giving identical bytes but different
content. Anyone adding a prefix must truncate at build time, or the retrievers
stop seeing the same string.

\subsection{For Evaluating Retrieval-Augmented Systems}

Three evaluation practices transfer. First, separate retrieval failure from
generation failure before scoring: most generations here fell in Case B, where
no gold anchor reached the top \thesisTopK{}, so a pooled ``hallucination rate''
would be mostly a retrieval statistic.

Second, annotate verbatim anchors, which, unlike chunk identifiers, stay stable
across chunk size (Methods), so one annotation scores all
\thesisNumConditions{} conditions.

Third, treat an automatic faithfulness proxy as an instrument with an error
rate that runs both ways: the census above reads faithful text under an
ungrounded label, and the misattribution reported in Results carries a grounded
and on-target one.

\section{Limitations and Future Research}

\subsection{Limitations}

The benchmark's rules select evidence: cleanly chunk-contained anchors, unique
values, prepared remarks over question-and-answer, management over analyst
speech, and, in eight items, question-to-passage overlap (Methods). Absolute retrieval scores are therefore likely
optimistic relative to a benchmark drawn without these rules; relative
comparisons between conditions remain informative, since every condition faced
the same questions. No claim is made about section effects. Two authoring gaps
were found after the benchmark was frozen, so neither can be repaired. The gold
anchor of \texttt{ref\_fact\_14} does not answer its own question, so the item
cannot score a correct answer as correct. It is one of the
\thesisNumQuestionsScored{} scored in every condition, and only I5 in
Table~\ref{tab:review-categories} turns on it alone.
And the difficulty annotation has a documented principle but no operational
rule, so it is usable as a label rather than a reproducible measure, and no analysis is keyed on it.
The first gap is a wrong value in one item; the second, a plausible value on
every item whose derivation cannot be audited.

The three hypothesis tests rest on only \thesisNumQuestionsScored{}
retrieval-scored questions, and in the reported arm many Stage 2 comparisons
were untestable outright, both conditions producing identical outcomes on every
matched question. The category breakdown's blocks differ in size,
\thesisNumQuestionsFactual{} factual,
\thesisNumQuestionsThematic{} thematic, and \thesisNumQuestionsComparative{}
comparative, and in achievable range: factual coverage@5 reached .786,
thematic never exceeded .133. The thematic null is a floor effect, and the
comparative result pointed the right way without surviving correction.

The Stage 1 test is itself contested in information retrieval evaluation.
\textcite{smucker2007} and \textcite{urbano2019} both conclude that the Wilcoxon
signed-rank test should be discontinued there and recommend randomization or
permutation tests instead, \textcite{smucker2007} because it is a simplified
variant of the randomization test. The objection is to the statistic, not to how
its \emph{p} value is computed, and every Stage 1 comparison here uses that
statistic, including every result reported as significant. \textcite{urbano2019}
keep the test legitimate where the paired differences are symmetric about zero,
which is the null under which \textcite{wilcoxon1945} introduced it, so the Stage
1 results stand only as far as that assumption holds for these data. The Stage 2
test is not exposed in the same way: on binary outcomes, the exact binomial form
used there coincides with a randomization test on the paired differences.

Grounding verifies that a number appears somewhere in the retrieved text, not
that it is attached to the right company, quarter, or metric in the answer's own
sentence. Per-claim attribution checking was out of scope, and the manual
review covers it on a subset only.

The results are conditioned on a single open-weights generator, and the
robustness check covers one open-weights model family and no other. Its second
model differs on several dimensions at once, so no difference can be attributed
to one of them.

The scope is narrow: \thesisNumCompanies{} United States banks over
\thesisNumQuarters{} quarters, one embedding model, one prompt, and a fixed
top~\thesisTopK{}, with a single annotator and no inter-annotator agreement. In
the enriched larger-chunk cell, build-time truncation reduced retrievability but
not scorability: anchor matching ran on raw text, never truncated.

\subsection{Future Research}

Exactly one enrichment format was tested, so which of its fields carry the
effect, whether a shorter prefix would retain it, and how prefix length trades
against the displacement described above remain unanswered. Varying the format
is the most direct follow-up.

The central claim is scoped to document sets that are alike. Running the same
grid over a corpus whose documents differ in subject matter would establish
where the effect disappears, and a null there would strengthen rather than
weaken the mechanism. The category split offers only a within-study hint at
that boundary, post hoc and confounded with a floor effect, which is precisely
why the between-corpus version is worth running. Cross-company retrieval remains
open beneath it: questions naming no company sat at floor in both arms, and the
comparative result pointed the right way without surviving correction, so what
would lift either is untested.

The census of ungrounded records shows how often a verbatim proxy misread
faithful text under that label in the reported arm. A faithfulness check that
tolerates paraphrase, normalization, and derivation, evaluated against that
census, would be the next instrument.

Conditioning Case C membership on the question's evidence requirement, rather
than on any single anchor, is a small and testable change to the scoring that
follows directly from the partial-context argument; testing it needs a
benchmark that records the requirement.
